\begin{trapbox}
  \textbf{\textcolor{CTrap}{易错 1｜只看 $\Delta<0$，忘了开口}}

  \textbf{错误想法：} $-x^2+2x-2$ 的 $\Delta=4-8=-4<0$，所以恒正。

  \textbf{实际：} $-x^2+2x-2=-(x-1)^2-1\le-1<0$，\textbf{恒负}！
  \begin{center}
  \begin{tikzpicture}[x=1cm,y=.63cm,>=stealth]
    \draw[->,gray] (-.65,0)--(2.75,0) node[right] {$x$};
    \draw[->,gray] (0,-3.5)--(0,.55) node[above] {$y$};
    \draw[red!75!black,very thick,domain=-.48:2.48,samples=45] plot (\x,{-(\x-1)*(\x-1)-1});
    \fill[red!75!black] (1,-1) circle (2pt);
    \node[anchor=west] at (3,-1.6) {最高也只有 $-1$};
  \end{tikzpicture}

  \end{center}
  \textbf{正确检查：} 恒正必须\textbf{$a>0$ 且 $\Delta<0$}。无交点不代表一定在轴上方。

  \medskip
  \textbf{\textcolor{CTrap}{易错 2｜把“$>0$”和“$\ge0$”混在一起}}

  \textbf{错误想法：} $\Delta\le0$ 就能保证恒正。

  \textbf{实际：} $x^2-2x+1=(x-1)^2$，当 $x=1$ 时恰好等于 $0$。
  \begin{center}
  \begin{tikzpicture}[x=1cm,y=.77cm,>=stealth]
    \draw[->,gray] (-.65,0)--(2.75,0) node[right] {$x$};
    \draw[->,gray] (0,-.5)--(0,2.75) node[above] {$y$};
    \draw[orange!85!black,very thick,domain=-.55:2.55,samples=45] plot (\x,{(\x-1)*(\x-1)});
    \fill[orange!85!black] (1,0) circle (2pt);
    \node[below] at (1,0) {$1$};
    \node[anchor=west] at (2.95,1.6) {$\Delta=0$：碰轴};
  \end{tikzpicture}

  \end{center}
  \textbf{正确检查：} 恒正用 $\Delta<0$；恒非负才允许 $\Delta=0$。

  \medskip
  \textbf{\textcolor{CTrap}{易错 3｜参数让二次项消失，还硬套判别式}}

  \textbf{原式：} $f(x)=(k-2)x^2+2x+1$。当 $k=2$ 时变为 $2x+1$，\textbf{不再是二次函数}。
  例如取 $x=-1$，就有 $2x+1=-1<0$。
  \begin{center}
  \begin{tikzpicture}[x=1.1cm,y=.57cm,>=stealth]
    \draw[->,gray] (-1.55,0)--(1.3,0) node[right] {$x$};
    \draw[->,gray] (0,-1.55)--(0,3.3) node[above] {$y$};
    \draw[red!75!black,very thick,domain=-1.2:1.1,samples=2] plot (\x,{2*\x+1});
    \fill[red!75!black] (-.5,0) circle (2pt);
    \node[below] at (-.5,0) {$-\frac12$};
    \node[anchor=west] at (1.5,1.45) {$k=2$：$y=2x+1$};
  \end{tikzpicture}

  \end{center}
  \textbf{正确检查：} 先分 $a=0$ 和 $a\ne0$。若 $a=0$，式子变成 $bx+c$：
  \textbf{$b\ne0$ 时不可能对全体实数恒正；$b=0,c>0$ 时反而恒正。}

  \medskip
  \textbf{\textcolor{CTrap}{易错 4｜题目只问一个区间，却套“全体实数”条件}}

  \textbf{题目换一换：} 判断 $x^2-1>0$ 对\textbf{$2\le x\le3$} 是否恒成立。
  虽然 $\Delta=4>0$，但在这个区间内 $x^2-1\ge3>0$，\textbf{确实恒成立}。
  \begin{center}
  \begin{tikzpicture}[x=.95cm,y=.34cm,>=stealth]
    \fill[green!15] (2,0) rectangle (3,9.2);
    \draw[->,gray] (-1.85,0)--(3.5,0) node[right] {$x$};
    \draw[->,gray] (0,-1.65)--(0,9.5) node[above] {$y$};
    \draw[blue!80!black,very thick,domain=-1.8:3.15,samples=65] plot (\x,{\x*\x-1});
    \draw[dashed,green!55!black] (2,0)--(2,3);
    \draw[dashed,green!55!black] (3,0)--(3,8);
    \fill[green!55!black] (2,3) circle (2.5pt);
    \fill[green!55!black] (3,8) circle (2.5pt);
    \node[below] at (2,0) {$2$};
    \node[below] at (3,0) {$3$};
    \node[anchor=west] at (3.6,6.6) {$[2,3]$ 上全在轴上方};
  \end{tikzpicture}

  \end{center}
  \textbf{正确检查：} 本结论只用于\textbf{所有实数}。若题目限制区间，要看该区间的图像，或比较端点和区间内顶点的函数值。
\end{trapbox}
